\chapter{Modal Verbs}

% dürfen (to be allowed to)
\begin{verbentry}{
  style = {modal},
  toc = {dürfen (to be allowed to)},
  name = {dürfen},
  english = {to be allowed to},
  type = {Modal},
  auxiliary = {haben}
}
 
     \vspace{5pt}
    \hrule
    \vspace{5pt}
  \begin{tabular}{rl}
     \multicolumn{2}{l}{\textbf{PRÄSENS} }\\
    ich & \textbf{darf}\\
    du & \textbf{darfst}\\
    er/sie/es & \textbf{darf}\\
    wir & \textbf{dürfen}\\
    ihr & \textbf{dürft}\\
    sie/Sie & \textbf{dürfen}\\
  \end{tabular}
  \hfill
     \begin{tabular}{rl}
    \multicolumn{2}{l}{\textbf{PRÄTERITUM} }\\
    ich & \textbf{durfte}\\
    du & \textbf{durftest}\\
    er/sie/es & \textbf{durfte}\\
    wir & \textbf{durften}\\
    ihr & \textbf{durftet}\\
    sie/Sie & \textbf{durften}\\
  \end{tabular}
   \hfill
     \begin{tabular}{rl}
    \multicolumn{2}{l}{\textbf{KONJUNKTIV II} }\\
    ich & \textbf{dürfte}\\
    du & \textbf{dürftest}\\
    er/sie/es & \textbf{dürfte}\\
    wir & \textbf{dürften}\\
    ihr & \textbf{dürftet}\\
    sie/Sie & \textbf{dürften}\\
  \end{tabular}
  \hfill
  \begin{tabular}{rl}
     \multicolumn{2}{l}{\textbf{IMPERATIV} }\\
   & \textbf{-}\\
   & \textbf{-}\\
   & \textbf{-}\\
   & \textbf{-}\\
   & \vspace{10pt}
  \end{tabular}

    \vspace{10pt}
 
    \begin{tabular}{rl}
     \multicolumn{2}{l}{\textbf{PERFEKT}}\\
   & haben + \textbf{gedurft}\\
  \end{tabular}
\hfill
   \begin{tabular}{rl}
   
    \multicolumn{2}{l}{\textbf{FUTURE I} }\\
    &werden + \textbf{dürfen}\\
 
  \end{tabular}
\hfill
   \begin{tabular}{rl}
      \multicolumn{2}{l}{\textbf{PLUSQUAMPERFEKT}}\\
   & hatten + \textbf{gedurft}\\
  \end{tabular}
  
  
\vspace{10pt}

\begin{tabular}{lll}
\multicolumn{3}{l}{\textbf{MIT PRÄPOSITIONEN}}\\
& \textbf{---} & \\
\end{tabular}
\hfill
\begin{tabular}{lll}
\multicolumn{3}{l}{\textbf{TRENNBARE VERBEN}}\\
& \textbf{---} & \\
\end{tabular}
\vspace{-5pt}
\end{verbentry}

% können (to be able to)
\begin{verbentry}{
  style = {modal},
  toc = {können (to be able to)},
  name = {können},
  english = {to be able to},
  type = {Modal},
  auxiliary = {haben}
}
 
     \vspace{5pt}
    \hrule
    \vspace{5pt}
  \begin{tabular}{rl}
     \multicolumn{2}{l}{\textbf{PRÄSENS} }\\
    ich & \textbf{kann}\\
    du & \textbf{kannst}\\
    er/sie/es & \textbf{kann}\\
    wir & \textbf{können}\\
    ihr & \textbf{könnt}\\
    sie/Sie & \textbf{können}\\
  \end{tabular}
  \hfill
     \begin{tabular}{rl}
    \multicolumn{2}{l}{\textbf{PRÄTERITUM} }\\
    ich & \textbf{konnte}\\
    du & \textbf{konntest}\\
    er/sie/es & \textbf{konnte}\\
    wir & \textbf{konnten}\\
    ihr & \textbf{konntet}\\
    sie/Sie & \textbf{konnten}\\
  \end{tabular}
   \hfill
     \begin{tabular}{rl}
    \multicolumn{2}{l}{\textbf{KONJUNKTIV II} }\\
    ich & \textbf{könnte}\\
    du & \textbf{könntest}\\
    er/sie/es & \textbf{könnte}\\
    wir & \textbf{könnten}\\
    ihr & \textbf{könntet}\\
    sie/Sie & \textbf{könnten}\\
  \end{tabular}
  \hfill
  \begin{tabular}{rl}
     \multicolumn{2}{l}{\textbf{IMPERATIV} }\\
   & \textbf{-}\\
   & \textbf{-}\\
   & \textbf{-}\\
   & \textbf{-}\\
   & \vspace{10pt}
  \end{tabular}

    \vspace{10pt}
 
    \begin{tabular}{rl}
     \multicolumn{2}{l}{\textbf{PERFEKT}}\\
   & haben + \textbf{gekonnt}\\
  \end{tabular}
\hfill
   \begin{tabular}{rl}
   
    \multicolumn{2}{l}{\textbf{FUTURE I} }\\
 &   werden + \textbf{können}\\
 
  \end{tabular}
\hfill
   \begin{tabular}{rl}
      \multicolumn{2}{l}{\textbf{PLUSQUAMPERFEKT}}\\
  &  hatten + \textbf{gekonnt}\\
  \end{tabular}
  
  
\vspace{10pt}

\begin{tabular}{lll}
\multicolumn{3}{l}{\textbf{MIT PRÄPOSITIONEN}}\\
& \textbf{---} & \\
\end{tabular}
\hfill
\begin{tabular}{lll}
\multicolumn{3}{l}{\textbf{TRENNBARE VERBEN}}\\
& \textbf{---} & \\
\end{tabular}
\vspace{-5pt}
\end{verbentry}

% mögen (to like)
\begin{verbentry}{
  style = {modal},
  toc = {mögen (to like)},
  name = {mögen},
  english = {to like},
  type = {Modal},
  auxiliary = {haben}
}
 
     \vspace{5pt}
    \hrule
    \vspace{5pt}
  \begin{tabular}{rl}
     \multicolumn{2}{l}{\textbf{PRÄSENS} }\\
    ich & \textbf{mag}\\
    du & \textbf{magst}\\
    er/sie/es & \textbf{mag}\\
    wir & \textbf{mögen}\\
    ihr & \textbf{mögt}\\
    sie/Sie & \textbf{mögen}\\
  \end{tabular}
  \hfill
     \begin{tabular}{rl}
    \multicolumn{2}{l}{\textbf{PRÄTERITUM} }\\
    ich & \textbf{mochte}\\
    du & \textbf{mochtest}\\
    er/sie/es & \textbf{mochte}\\
    wir & \textbf{mochten}\\
    ihr & \textbf{mochtet}\\
    sie/Sie & \textbf{mochten}\\
  \end{tabular}
    \hfill
     \begin{tabular}{rl}
    \multicolumn{2}{l}{\textbf{KONJUNKTIV II} }\\
    ich & \textbf{möchte}\\
    du & \textbf{möchtest}\\
    er/sie/es & \textbf{möchte}\\
    wir & \textbf{möchten}\\
    ihr & \textbf{möchtet}\\
    sie/Sie & \textbf{möchten}\\
  \end{tabular}
  \hfill
  \begin{tabular}{rl}
     \multicolumn{2}{l}{\textbf{IMPERATIV} }\\
   & \textbf{-}\\
   & \textbf{-}\\
   & \textbf{-}\\
   & \textbf{-}\\
   & \vspace{10pt}
  \end{tabular}

    \vspace{10pt}
 
    \begin{tabular}{rl}
     \multicolumn{2}{l}{\textbf{PERFEKT}}\\
  &  haben + \textbf{gemocht}\\
  \end{tabular}
\hfill
   \begin{tabular}{rl}
   
    \multicolumn{2}{l}{\textbf{FUTURE I} }\\
  &  werden + \textbf{mögen}\\
 
  \end{tabular}
\hfill
   \begin{tabular}{rl}
      \multicolumn{2}{l}{\textbf{PLUSQUAMPERFEKT}}\\
   & hatten + \textbf{gemocht}\\
  \end{tabular}
  
  
\vspace{10pt}

\begin{tabular}{lll}
\multicolumn{3}{l}{\textbf{MIT PRÄPOSITIONEN}}\\
& \textbf{---} & \\
\end{tabular}
\hfill
\begin{tabular}{lll}
\multicolumn{3}{l}{\textbf{TRENNBARE VERBEN}}\\
& \textbf{---} & \\
\end{tabular}
\vspace{-5pt}
\end{verbentry}

% müssen (to have to)
\begin{verbentry}{
  style = {modal},
  toc = {müssen (to have to)},
  name = {müssen},
  english = {to have to},
  type = {Modal},
  auxiliary = {haben}
}
 
     \vspace{5pt}
    \hrule
    \vspace{5pt}
  \begin{tabular}{rl}
     \multicolumn{2}{l}{\textbf{PRÄSENS} }\\
    ich & \textbf{muss}\\
    du & \textbf{musst}\\
    er/sie/es & \textbf{muss}\\
    wir & \textbf{müssen}\\
    ihr & \textbf{müsst}\\
    sie/Sie & \textbf{müssen}\\
  \end{tabular}
  \hfill
     \begin{tabular}{rl}
    \multicolumn{2}{l}{\textbf{PRÄTERITUM} }\\
    ich & \textbf{musste}\\
    du & \textbf{musstest}\\
    er/sie/es & \textbf{musste}\\
    wir & \textbf{mussten}\\
    ihr & \textbf{musstet}\\
    sie/Sie & \textbf{mussten}\\
  \end{tabular}
   \hfill
     \begin{tabular}{rl}
    \multicolumn{2}{l}{\textbf{KONJUNKTIV II} }\\
    ich & \textbf{müsste}\\
    du & \textbf{müsstest}\\
    er/sie/es & \textbf{müsste}\\
    wir & \textbf{müssten}\\
    ihr & \textbf{müsstet}\\
    sie/Sie & \textbf{müssten}\\
  \end{tabular}
  \hfill
  \begin{tabular}{rl}
     \multicolumn{2}{l}{\textbf{IMPERATIV} }\\
   & \textbf{-}\\
   & \textbf{-}\\
   & \textbf{-}\\
   & \textbf{-}\\
   & \vspace{10pt}
  \end{tabular}

    \vspace{10pt}
 
    \begin{tabular}{rl}
     \multicolumn{2}{l}{\textbf{PERFEKT}}\\
  &  haben + \textbf{gemusst}\\
  \end{tabular}
\hfill
   \begin{tabular}{rl}
   
    \multicolumn{2}{l}{\textbf{FUTURE I} }\\
  &  werden + \textbf{müssen}\\
 
  \end{tabular}
\hfill
   \begin{tabular}{rl}
      \multicolumn{2}{l}{\textbf{PLUSQUAMPERFEKT}}\\
  &  hatten + \textbf{gemusst}\\
  \end{tabular}
  
  
\vspace{10pt}

\begin{tabular}{lll}
\multicolumn{3}{l}{\textbf{MIT PRÄPOSITIONEN}}\\
& \textbf{---} & \\
\end{tabular}
\hfill
\begin{tabular}{lll}
\multicolumn{3}{l}{\textbf{TRENNBARE VERBEN}}\\
& \textbf{---} & \\
\end{tabular}
\vspace{-5pt}
\end{verbentry}

% sollen (should)
\begin{verbentry}{
  style = {modal},
  toc = {sollen (should)},
  name = {sollen},
  english = {should},
  type = {Modal},
  auxiliary = {haben}
}
 
     \vspace{5pt}
    \hrule
    \vspace{5pt}
  \begin{tabular}{rl}
     \multicolumn{2}{l}{\textbf{PRÄSENS} }\\
    ich & \textbf{soll}\\
    du & \textbf{sollst}\\
    er/sie/es & \textbf{soll}\\
    wir & \textbf{sollen}\\
    ihr & \textbf{sollt}\\
    sie/Sie & \textbf{sollen}\\
  \end{tabular}
  \hfill
     \begin{tabular}{rl}
    \multicolumn{2}{l}{\textbf{PRÄTERITUM} }\\
    ich & \textbf{sollte}\\
    du & \textbf{solltest}\\
    er/sie/es & \textbf{sollte}\\
    wir & \textbf{sollten}\\
    ihr & \textbf{solltet}\\
    sie/Sie & \textbf{sollten}\\
  \end{tabular}
   \hfill
     \begin{tabular}{rl}
    \multicolumn{2}{l}{\textbf{KONJUNKTIV II} }\\
    ich & \textbf{sollte}\\
    du & \textbf{solltest}\\
    er/sie/es & \textbf{sollte}\\
    wir & \textbf{sollten}\\
    ihr & \textbf{solltet}\\
    sie/Sie & \textbf{sollten}\\
  \end{tabular}
  \hfill
  \begin{tabular}{rl}
     \multicolumn{2}{l}{\textbf{IMPERATIV} }\\
   & \textbf{soll(e) du!}\\
   & \textbf{sollen wir!}\\
   & \textbf{sollt (ihr)}\\
   & \textbf{sollen Sie!}\\
   & \vspace{10pt}
  \end{tabular}

    \vspace{10pt}
 
    \begin{tabular}{rl}
     \multicolumn{2}{l}{\textbf{PERFEKT}}\\
  &  haben + \textbf{gesollt}\\
  \end{tabular}
\hfill
   \begin{tabular}{rl}
   
    \multicolumn{2}{l}{\textbf{FUTURE I} }\\
  &  werden + \textbf{sollen}\\
 
  \end{tabular}
\hfill
   \begin{tabular}{rl}
      \multicolumn{2}{l}{\textbf{PLUSQUAMPERFEKT}}\\
  &  hatten + \textbf{gesollt}\\
  \end{tabular}
  
  
\vspace{10pt}

\begin{tabular}{lll}
\multicolumn{3}{l}{\textbf{MIT PRÄPOSITIONEN}}\\
& \textbf{---} & \\
\end{tabular}
\hfill
\begin{tabular}{lll}
\multicolumn{3}{l}{\textbf{TRENNBARE VERBEN}}\\
& \textbf{---} & \\
\end{tabular}
\vspace{-5pt}
\end{verbentry}

% wollen (to want)
\begin{verbentry}{
  style = {modal},
  toc = {wollen (to want)},
  name = {wollen},
  english = {to want},
  type = {Modal},
  auxiliary = {haben}
}
 
     \vspace{5pt}
    \hrule
    \vspace{5pt}
  \begin{tabular}{rl}
     \multicolumn{2}{l}{\textbf{PRÄSENS} }\\
    ich & \textbf{will}\\
    du & \textbf{willst}\\
    er/sie/es & \textbf{will}\\
    wir & \textbf{wollen}\\
    ihr & \textbf{wollt}\\
    sie/Sie & \textbf{wollen}\\
  \end{tabular}
  \hfill
     \begin{tabular}{rl}
    \multicolumn{2}{l}{\textbf{PRÄTERITUM} }\\
    ich & \textbf{wollte}\\
    du & \textbf{wolltest}\\
    er/sie/es & \textbf{wollte}\\
    wir & \textbf{wollten}\\
    ihr & \textbf{wolltet}\\
    sie/Sie & \textbf{wollten}\\
  \end{tabular}
   \hfill
     \begin{tabular}{rl}
    \multicolumn{2}{l}{\textbf{KONJUNKTIV II} }\\
    ich & \textbf{wollte}\\
    du & \textbf{wolltest}\\
    er/sie/es & \textbf{wollte}\\
    wir & \textbf{wollten}\\
    ihr & \textbf{wolltet}\\
    sie/Sie & \textbf{wollten}\\
  \end{tabular}
  \hfill
  \begin{tabular}{rl}
     \multicolumn{2}{l}{\textbf{IMPERATIV} }\\
   & \textbf{wolle! (du)}\\
   & \textbf{wollen wir!}\\
   & \textbf{wollt! (ihr)}\\
   & \textbf{wollen Sie!}\\
   & \vspace{10pt}
  \end{tabular}

    \vspace{10pt}
 
    \begin{tabular}{rl}
     \multicolumn{2}{l}{\textbf{PERFEKT}}\\
  &  haben + \textbf{gewollt}\\
  \end{tabular}
\hfill
   \begin{tabular}{rl}
   
    \multicolumn{2}{l}{\textbf{FUTURE I} }\\
  &  werden + \textbf{wollen}\\
 
  \end{tabular}
\hfill
   \begin{tabular}{rl}
      \multicolumn{2}{l}{\textbf{PLUSQUAMPERFEKT}}\\
  &  hatten + \textbf{gewollt}\\
  \end{tabular}
  
  
\vspace{10pt}

\begin{tabular}{lll}
\multicolumn{3}{l}{\textbf{MIT PRÄPOSITIONEN}}\\
& \textbf{---} & \\
\end{tabular}
\hfill
\begin{tabular}{lll}
\multicolumn{3}{l}{\textbf{TRENNBARE VERBEN}}\\
& \textbf{---} & \\
\end{tabular}
\vspace{-5pt}
\end{verbentry}
